\section{Introduction}\label{sec:intro}

Estimating a linear function of the state, rather than the whole state, is often what a
control loop actually needs. Functional observers do exactly this, and their appeal is
quantitative: the order of a functional observer can be far below the order of a
Luenberger observer for the same system, because only the directions that the functional
sees have to be estimated. The design question is therefore an order question: given the
system and the functionals to be estimated, what is the smallest number of observer rows
for which a functional observer exists?

For a linear system with measured output, Fernando, Trinh and Jennings \cite{fer10}
reduced that question to two rank conditions on the observer matrix: with the observer
rows written as a block $L$ that contains the given functional rows $L_0$, one condition
is algebraic and one is spectral. They characterize the feasible rows for a prescribed
order, but they do not address the cost of deciding whether such rows exist.

The characterization itself has since been completed. A unified algebraic framework is
given in \cite{dar25}. The two recent works \cite{fer26a,fer26b} determine, order by
order, which orders are attainable: \cite{fer26a} identifies through functional
observability indices the minimum dimension permitted by the algebraic condition, the
complete family of augmentations attaining it, and necessary and sufficient spectral
conditions for that dimension to be optimal, while \cite{fer26b} characterizes the
interval between that minimum and the endpoint set by functional observability, using a
nullspace representation of the algebraic condition and a matrix-pencil form of the
spectral one. The latter work also describes an order-by-order procedure: at each order
it parameterizes the nullspace over Grassmann charts with $d(m_L-d)=(n-p-q)(q-r)$ free
parameters per chart and tests the two conditions chart by chart.

Neither work states a running-time or bit-complexity bound, and for fixed $q$ the number
of chart parameters above still grows linearly in $n$. The complexity of deciding
feasibility at a prescribed order is therefore left open by the characterization
literature. This paper takes the conditions as given and studies that complexity.

Adjacent design problems are known to be hard, which makes the question natural but does
not answer it. Choosing a minimal set of sensors that renders a system functionally
observable, or structurally functionally observable, is NP-hard \cite{zha25}, and
constrained sensor-placement variants are NP-complete \cite{dey21}. Static output
feedback with a spectral predicate is NP-hard \cite{lof26,bar26} and, in its real
feasibility form, $\exists\mathbb R$-complete \cite{liu26}; affine rank minimization is
$\exists\mathbb R$-complete already at low rank \cite{maj26}. None of these results
transfers. In sensor placement the quantified object is a subset of the rows of the
output matrix, so the output matrix is itself a design variable, whereas here it is fixed
and the design variable is the set of observer rows. The predicates also differ: a
condition on the spectral abscissa of a closed loop, against a constant-rank condition
that has to hold for every complex value of the spectral variable. A reduction would
moreover have to preserve the single block of free rows that occurs in both conditions,
together with the common origin of the compressed matrices, whereas rank-constrained
formulations leave every entry of the constrained matrix free.

The approach is to make the decision problem finite-dimensional. Restricting row
functionals to the common kernel of the output matrix and its product with the system
matrix removes exactly the directions that those two matrices already account for, and
turns the two conditions into a row-space inclusion and into the constant rank of a
pencil. Every free observer row is then determined by its coordinates in the space of
attainable row pairs, whose dimension gives the effective dimension $\kappa_{\rm obs}$ of
the free data (Section~\ref{sec:dimension}). Two consequences are developed: the spectral
condition becomes a finite algebraic certificate compressed to a size governed by
$\kappa_{\rm obs}$ (Section~\ref{sec:certificate}), and deciding feasibility becomes a
real-algebraic problem whose number of variables is controlled by either the prescribed
order or the effective dimension (Sections~\ref{sec:fixedq} and~\ref{sec:fpt}).

\subsection{Contributions and organization}\label{sec:contrib}

The contributions are the following.

\begin{enumerate}
\item The identification of the effective dimension $\kappa_{\rm obs}$ of the free
observer data as the structural parameter of the problem, and a compression of the two
observer conditions, which are taken from the existing characterization, to $O(q^2)$
variables, independently of the state dimension, which the chart parameterization of
\cite{fer26b} does not achieve.
\item A finite algebraic certificate for the spectral condition, and its compression by
Cauchy--Binet to a matrix of size controlled by $\kappa_{\rm obs}$; no enumeration of
maximal minors is required.
\item A decision algorithm for every fixed order, with a bit-complexity bound
$\mathcal L^{O(q^2)}$ in the input length $\mathcal L$ and the prescribed order $q$.
\item A decision algorithm that is fixed-parameter tractable in $\kappa_{\rm obs}$, with
bound $f(\kappa_{\rm obs})\operatorname{poly}(\mathcal L)$, together with an exact
polynomial-time test for the case of no free rows, which decides a finite rational linear
system obtained from a polynomial left inverse.
\item Membership of the prescribed-order decision problem, and of its finite threshold
version, in the existential theory of the reals.
\end{enumerate}

The paper is organized as follows. Section~\ref{sec:prelim} fixes the notation and states
the decision, optimization and construction problems. Section~\ref{sec:compressed} proves
the compression, Section~\ref{sec:dimension} identifies the effective dimension, and
Section~\ref{sec:certificate} develops the certificate. Sections~\ref{sec:fixedq}
and~\ref{sec:fpt} give the two algorithms, Section~\ref{sec:existsr} records the
existential-real membership, and Section~\ref{sec:conclusions} concludes.

\section{Preliminaries}\label{sec:prelim}

\subsection{Notation and conventions}\label{sec:notation}

Consider the linear time-invariant system
\begin{equation}\label{system}
\dot{x}=Ax+Bu,\quad y=Cx,\quad z=L_0x,
\end{equation}
where $A\in\mathbb Q^{n\times n}$, $C\in\mathbb Q^{p\times n}$ and
$L_0\in\mathbb Q^{r_0\times n}$.

Let $q\ge r_0$ be the prescribed observer order, with $q\le n$, and set $r=q-r_0$. The
observer rows are the rows of the block
\begin{equation*}
L=[L_0;\,R],\qquad R\in\mathbb R^{r\times n},
\end{equation*}
where $L_0$ is fixed and $R$ collects the $r$ free rows; only real observers are
considered.

All given rational data are encoded in binary, and $\mathcal L$ denotes the total input
bit length. The free parameter $R$ is not part of the given input data and is therefore
not counted in $\mathcal L$. State vectors are column vectors, and matrices act on them by
the usual left multiplication.

Let $P(\lambda)$ be a rational matrix pencil of generic rank $\omega$, and call its
$\omega\times\omega$ minors maximal minors. Ranks at spectral points $\lambda\in\mathbb C$
are always computed over $\mathbb C$. Because the maximal minors have rational
coefficients, their gcd is a nonzero constant in $\mathbb C[\lambda]$ exactly when it is a
nonzero constant in $\mathbb Q[\lambda]$; this equivalence is used without further
mention.

\subsection{The observer conditions}\label{sec:conditions}

A functional observer of order $q$ is a matrix $L=[L_0;R]\in\mathbb R^{q\times n}$ such
that
\begin{equation}\label{eq:5a}
\operatorname{rank}\begin{bmatrix}LA\\ L\\ CA\\ C\end{bmatrix}
=\operatorname{rank}\begin{bmatrix}L\\ CA\\ C\end{bmatrix},
\end{equation}
\begin{equation}\label{eq:5b}
\operatorname{rank}\begin{bmatrix}\lambda L-LA\\ CA\\ C\end{bmatrix}
=\operatorname{rank}\begin{bmatrix}L\\ CA\\ C\end{bmatrix},
\quad \forall\lambda\in\mathbb C.
\end{equation}
These are the Fernando--Trinh--Jennings conditions, numbered (5a) and (5b) in
\cite{fer10}. Condition~\eqref{eq:5a} is algebraic; \eqref{eq:5b} is spectral, and it is
required at every $\lambda\in\mathbb C$, in particular at all finite eigenvalues of $A$.

\subsection{Decision, optimization and construction}\label{sec:problems}

Three problems are posed for \eqref{eq:5a} and \eqref{eq:5b}, in increasing order of
demand.

\emph{Decision.} $\mathrm{DEC}$ takes $(A,C,L_0)$ and $q$ and asks whether some $R$
satisfies \eqref{eq:5a} and \eqref{eq:5b}; no witness is required. This is the problem
studied in this paper.

\emph{Optimization.} $\mathrm{OPT}$ asks for the minimum feasible order
\begin{equation*}
q_{\min}=\min\{q:\mathrm{DEC}(q)\},
\end{equation*}
where $\mathrm{DEC}(q)$ abbreviates $\mathrm{DEC}(A,C,L_0,q)$.

\emph{Construction.} $\mathrm{FIND}$ asks for a witness $R$ at the minimum feasible
order, $q=q_{\min}$.

The threshold problem is
\begin{equation}\label{eq:mindec}
\mathrm{MINDEC}(Q)=\bigvee_{q=r_0}^{\min(Q,n)}\mathrm{DEC}(q),
\end{equation}
a finite disjunction of decision instances. The algorithms below decide $\mathrm{DEC}$;
their affirmative certificates contain a witness. We make no claim here about the
complexity of $\mathrm{OPT}$ or $\mathrm{FIND}$, nor do we rely on monotonicity of
$\mathrm{DEC}$ in $q$.
